% ============================================================================
%  2. METHODS
%
%  Reglas A&A sobre ecuaciones:
%    - Toda ecuación referenciada con \ref necesita su \label. Solo ese
%      mecanismo; nada de numerar a mano.
%    - Puntuar las ecuaciones como si fueran texto: coma o punto al final.
%    - Delimitadores con \left( ... \right) para que escalen con la fórmula.
%    - En modo matemático, todo va en itálica por defecto. Lo que NO es una
%      variable va en redonda:
%          \mathrm{d}   diferencial
%          \mathrm{e}   base del logaritmo natural
%          \mathrm{i}   unidad imaginaria
%          $T_\mathrm{eff}$, $m_\mathrm{e}$   subíndices que son etiquetas
%          \mbox{Hz}, \mathrm{km\,s^{-1}}     unidades
%      Pero SÍ en itálica si el subíndice es a su vez una variable.
%    - Vectores: \vec{A}.  Tensores: \tens{ABC}.  Solo en modo matemático.
%    - La numeración de ecuaciones en A&A es correlativa en todo el artículo,
%      no por sección.
% ============================================================================

\section{Methods}
\label{sec:methods}

Texto de la sección.

% ----------------------------------------------------------------------------
%  Ejemplo de ecuación numerada y referenciable
% ----------------------------------------------------------------------------
% \begin{equation}
%   \label{eq:geodesic}
%   \frac{\mathrm{d}^2 x^\mu}{\mathrm{d}\tau^2}
%     + \Gamma^\mu_{\ \alpha\beta}
%       \frac{\mathrm{d}x^\alpha}{\mathrm{d}\tau}
%       \frac{\mathrm{d}x^\beta}{\mathrm{d}\tau} = 0 \,,
% \end{equation}
%
% donde $\tau$ es el tiempo propio.

\subsection{Subsección}
\label{sec:methods-sub}

Texto.
